% ============================================================
%  sec01.tex —— §1 引言（原文 pp. 2–6，止于第 2 节标题之前）
%  原文：Homological Mirror Symmetry for Elliptic Curves（2013-04-25）
%  译者：t1-intro        公式已逐条对照 p02–p06 页面图像核对
% ============================================================

\section{引言}\label{sec:intro}

\begin{keytakeaways}
  \item \textbf{物理/几何图像}：镜像对称把 $X$ 上的代数几何（B-模型）与镜像辛流形 $X^\vee$ 上的辛几何（A-模型）对应起来；式 \eqref{eq:intro-1} 只是这一对应的口号，Kontsevich 的纲领是把它升级为两个范畴之间的等价。
  \item \textbf{数学工具}：代数侧取凝聚层的有界导出范畴 $D^b(Coh(X))$（三角范畴），辛侧取分裂闭包的三角导出 Fukaya 范畴 $H^0D^\pi\Fuk(X^\vee)$；链级陈述则改用 dg-范畴 $D_\infty(X)$ 与 $A_\infty$-范畴，上同调层面的结果只是链级现象的“影子”。
  \item \textbf{核心结论}：在椭圆曲线情形可以证明链级镜像等价 $D_\infty(X)\to D^\pi\Fuk(X^\vee)$；证明的枢纽是两侧范畴都由同一个有限维箭图代数 $Q$ 生成，问题化为比较 $Q$ 上的两个 $A_\infty$-结构。
  \item \textbf{关键机制}：先用同伦微扰论把链级信息压缩成 $Q$ 上的高阶乘法 $(m_d)$，再用 Hochschild 上同调对 $A_\infty$-结构分类，最后用 Lekili–Perutz 的“几乎所有 $A_\infty$-结构都来自椭圆曲线”这一定理闭合论证。
\end{keytakeaways}

镜像对称最早是由研究所谓超共形弦论（superconformal string theory）的物理学家注意到的。特别地，给定一个代数卡拉比–丘流形（algebraic Calabi–Yau manifold）$X$，物理学家预言存在一个“镜像”辛流形 $X^\vee$（它也是卡拉比–丘流形），以及两者之间的种种关系；这些关系可以粗略地概括为
\begin{equation}\label{eq:intro-1}
  \text{$X$ 上的代数几何}\ \longleftrightarrow\ \text{$X^\vee$ 上的辛几何}.
\end{equation}

% [Tutor Note: 式 \eqref{eq:intro-1} 的 $\longleftrightarrow$ 只是口号，不能读成某个具体函子：左端是 $X$ 上的复几何（B-模型），右端是 $X^\vee$ 上的辛几何（A-模型），两侧方向不可交换。另注意 $\vee$ 在本文有两种用法——镜像流形 $X^\vee$ 与“对偶”——直到“记号与约定”才统一交代；$X^\vee$ 并非线性对偶或层论对偶，读后文时不要混淆。]

% [Tutor Note: 给物理背景读者的词典：$D^b(Coh(X))$ 的对象是 B-型 D-膜（$X$ 的复子流形上的凝聚层，向量丛是其中一类），$\Fuk(X^\vee)$ 的对象是 A-型 D-膜（$X^\vee$ 的拉格朗日子流形 $L$，连同其上的平坦线丛）。条件 $\omega_Y|_L=0$ 就是“$L$ 是拉格朗日”的定义式，也是 A-膜能存在于辛流形中的几何前提；HMS 断言这两类 D-膜的范畴等价，而 \eqref{eq:intro-1} 只是它在“几何”层面的投影。]

文献 [3] 是一本试图补足上述粗略概括中“$\leftrightarrow$”的含糊之处的书。

Kontsevich 在 1994 年国际数学家大会（ICM）上给出了（表面上）颇为不同的一条澄清途径 [10]。他提出一组猜想，如今统称为“同调镜像对称”（homological mirror symmetry）：这些猜想把物理学家的许多论断提炼进一个范畴框架。下面简述之。

我们从代数卡拉比–丘簇 $X$ 出发。编码 $X$ 的几乎全部几何信息的一个范畴，是其凝聚层（coherent sheaf）的有界导出范畴（bounded derived category）$D^b(Coh(X))$。它是在 $X$ 上凝聚层的有界上链复形范畴中把拟同构（quasi-isomorphism）反演成同构而得到的范畴。这不是阿贝尔范畴；但它是加性范畴，并带有某种额外结构，即平移函子（shift functor）与特殊三角形（distinguished triangle）。这些额外数据通常在“三角范畴”（triangulated category）这一名目下公理化。（[23] 第十章是理解这个意义上的三角范畴与导出范畴的好参考，但要留意：本文正文中我们通常另有所指。）这就是 Kontsevich 对“式”\eqref{eq:intro-1} 右端的解释：把 $X$ 上凝聚层的有界导出范畴看作三角范畴。

% [Tutor Note: 原文疑误：本段称 $D^b(Coh(X))$ 是 (1) 的“右端”，但 (1) 的排版是“$X$ 上代数几何 $\longleftrightarrow$ $X^\vee$ 上辛几何”，代数侧在左；下一段又把 $\Fuk$ 范畴称为 (1) 的“右端”。两处必有一误，按上下文此处应为“左端”。译文保留原文并标注。]

% [Tutor Note: “把拟同构反演成同构”即 Verdier 局部化：$D^b(Coh(X))$ 由 $K^b(Coh(X))$ 沿拟同构构成的态射系作局部化得到，三角结构（平移函子 + 特殊三角形）正是从复形范畴继承下来的。作者在括号里特意提醒 [23] 的“三角范畴”含义与正文不同：正文后面要用 dg/$A_\infty$-增强取代抽象三角范畴，读 §3 时须记住这条预警。]

现在考虑一个辛卡拉比–丘流形 $Y$：即一个第一陈类为零的实光滑流形，连同所选定的一个闭的非退化二形式 $\omega_Y$。

% [Tutor Note: 这里的 $\omega_Y$ 是复化后的辛形式，按约定它与近复结构相容（后文 §4 才细说）；$Y$ 的“卡拉比–丘”条件（$c_1=0$）保证 Maslov 类可定义，从而拉格朗日子流形可以分次化——这是 $\Fuk(Y)$ 能赋 $\ZZ$-分次、进而与 $D^b(Coh(X))$ 的分次结构比较的前提。]

在 Fukaya 工作的基础上，Kontsevich 猜想存在所谓分裂闭包的三角导出 Fukaya 范畴（split-closed triangulated derived Fukaya category）$H^0D^\pi\Fuk(Y)$；这就是 Kontsevich 对 \eqref{eq:intro-1} 右端的解释。只要弄清楚 $H^0D^\pi\Fuk(Y)$ 是什么，同调镜像对称猜想就可以表述如下：

\begin{conjecture*}[Kontsevich]
设 $X$ 为射影卡拉比–丘流形。则存在辛流形 $X^\vee$，以及三角范畴的等价
\[
  D^b(Coh(X)) \simeq H^0D^\pi\Fuk(X^\vee).
\]
\end{conjecture*}

% [Tutor Note: $H^0D^\pi\Fuk(Y)$ 要逐层读：$\Fuk(Y)$ 是链级 $A_\infty$-范畴；$D^\pi$ 是“分裂闭包的三角化包络”（$\pi$ 指幂等分裂，即把幂等态射的形式像添为新对象）；最后 $H^0$ 取第 0 次上同调，才得到普通的三角范畴。三步缺一不可：若只取 $H^0\Fuk(Y)$ 而不做 $D^\pi$，所得范畴既不幂等完备也不三角。]

% [Tutor Note: 易错点（方向）：本猜想把代数侧写在左边、辛侧写在右边，即函子方向为 $D^b(Coh(X))\to H^0D^\pi\Fuk(X^\vee)$。文献中有作者把等价反向写成 $\Fuk(X^\vee)\to D^b(Coh(X))$，甚至在 A 侧取 opposite category（因 Floer 复合的书写次序与层论复合相反）。引用或比较不同文献的 HMS 陈述时，务必先对齐范畴与函子的方向。]

作为一阶近似，$\Fuk(Y)$ 的对象由 $Y$ 的紧拉格朗日子流形（Lagrangian submanifold）组成（即满足 $\omega_Y|_L=0$ 的子流形 $L$）。特别地，它们是半维的；因此在好的情形（横截相交）下，两个拉格朗日子流形 $L_1,L_2$ 的交点集合 $L_1\cap L_2$ 是有限的。我们把这些交点取作某个向量空间的形式生成元；这个向量空间通常称为 Floer 链复形（Floer chain complex），记作 $\Hom_{\Fuk(Y)}(L_1,L_2)$。顾名思义，它带有微分，并且将有
\[
  \Hom_{H^0D^\pi\Fuk(Y)}(L_1,L_2) = H^0\Hom_{\Fuk(Y)}(L_1,L_2).
\]
也就是说，Kontsevich 猜想中的那个导出 Fukaya 范畴，是把“链级”范畴 $\Fuk(Y)$ 取上同调得到的。

% [Tutor Note: 易错点：$\Hom_{\Fuk(Y)}(L_1,L_2)$ 是复形（带微分、带 $\ZZ$-分次），而“记号与约定”里所说的 $\Hom$ 指分次向量空间的 0 次映射；同一个符号在链级与上同调级含义不同。上式只是把“链级 $\to$ 上同调级”这句话写成公式。具体地，微分计数全纯条带（bigon），复合计数三角形，即原文所谓的 immersed bigons/triangles；在椭圆曲线情形这些计数退化为组合的多边形面积计数。]

以上做法遮掩了种种困难，我们至少要把它们提一下。(1) 当 $L_1$ 与 $L_2$ 不横截相交时存在明显的问题，而这实际上极难修补。(2) 我们还没有定义微分，也没有定义态射的复合。在好的情形，这可以由拉格朗日相交 Floer 理论（Lagrangian intersection Floer theory）做到：微分与复合通过计数浸入的二边形与三角形（在某种等价关系下）来计算。取上同调之后，这确实给出一个通常的范畴；但链级结构非常复杂，特别是链级复合不满足结合律。补救的办法是引入更高阶的复合运算（同样由计数边数更多的多边形来定义），它们充当度量结合律失效程度的同伦。由此得到的代数结构通常称为 $A_\infty$-范畴（$A_\infty$-category）。(3) 事实证明，范畴 $H^0\Fuk(Y)$ 并不特别好：例如幂等（idempotent）自同态不必对应子对象；又如 $H^0\Fuk(Y)$ 虽然是加性范畴，却既不是阿贝尔的，也不是三角的。不过存在一个纯代数的手续，称为“取分裂闭包的三角化包络”（forming the split-closed triangulated envelope），它从 $\Fuk(Y)$ 构造出新的 $A_\infty$-范畴 $D^\pi\Fuk(Y)$，使得 $H^0D^\pi\Fuk(Y)$ 具有上述良好性质（即带有自然的三角结构，且幂等自同态对应子对象）。正是这最后一个范畴，才被期望与有界导出范畴三角等价。

% [Tutor Note: 作者说 $D^\pi\Fuk(Y)$“具有良好性质”时，对象仍是 $A_\infty$-范畴，良好性质属于 $H^0D^\pi\Fuk(Y)$：三角化包络并不是把 $A_\infty$-范畴直接变成三角范畴，而是添入锥（twisted complexes）与幂等像，只有在上同调层面才呈现三角结构。严格构造见 Seidel, \textit{Fukaya categories and Picard--Lefschetz theory} 中关于 twisted complexes 与 split-closure 的章节。]

按目前的表述，Kontsevich 的猜想有一处缺陷。这与下述观察有关：上同调层面的结果，说到底应当只是链级现象的“影子”。举例来说，导出范畴虽然是一个陈述结果的自然的范畴，用起来其实相当别扭（例如 $X$ 的导出范畴无法由某个开覆盖中各块的导出范畴“粘”回来）。因此应当设法找到一个自然的“链级”范畴 $D_\infty(X)$，使得
\[
  H^0D_\infty(X) \simeq D^b(Coh(X)).
\]

% [Tutor Note: “不能用开覆盖粘回来”指 $D^b(Coh(X))$ 不满足 Zariski/étale 下降；要使下降成立必须使用 enhanced/derived 结构（Lurie、Toën 的高结构）。这是引入链级范畴 $D_\infty(X)$ 的实质动机之一，而不只是为了“让定理的表述更漂亮”。]

幸运的是，这样的范畴是众所周知的（事实上它们正是最初用来证明 $D^b(Coh(X))$ 存在性的工具）。此外，$A_\infty$-范畴的概念显然适用于 $D_\infty(X)$（我们实际上得到一种特别简单的 $A_\infty$-结构，称为 dg-范畴（dg-category），即微分分次范畴）。于是可以把 Kontsevich 的猜想改进为如下形式：

\begin{conjecture*}[Kontsevich，链级]
设 $X$ 为射影卡拉比–丘流形。则存在辛流形 $X^\vee$，以及 $A_\infty$-范畴的态射
\[
  D_\infty(X) \to D^\pi\Fuk(X^\vee),
\]
它在取上同调之后诱导同构。
\end{conjecture*}

这个猜想如今已对（射影空间中的）所有卡拉比–丘超曲面得到验证。椭圆曲线（维数一）的情形知道得最早，由 Polishchuk 与 Zaslow 的一系列论文 [17][14][15][16] 建立。四次三维簇（维数二）的情形由 Seidel [19] 处理。维数 $\ge 3$ 的情形由 Sheridan [21] 建立。

% [Tutor Note: 原文疑误：Seidel [19] 处理的是四次曲面（quartic surface，即复二维的 K3），而不是“四次三维簇”；本段“Quartic threefolds (dimension two)”自相矛盾（三次簇是复三维），且后文第 4 页自称 “Seidel for quartic surfaces”。译文保留原文并标注。同段 “in a series papers” 应为 “in a series of papers”，系笔误。]

本文的主要目标，是给出 Kontsevich 猜想第一个非平凡情形（即椭圆曲线）的证明。做法是把 Seidel 为四次曲面（quartic surface）发展的精细方法，应用到远比它容易的一维情形上。在此过程中我们基本上遵循 [11]，至少在精神上如此（实践上我们只打算证明少得多的东西，因而可以省去一些枝节）。

本文还有两个次要目标。第一，铺陈出为精确表述并最终证明 Kontsevich 猜想所需的、相对大量的背景材料。第二，我们想同时说明两件事：同调镜像对称的证明中所使用的技巧，以及为辨认镜像对两侧的辛几何与代数几何所需要的那种非同寻常的“巧合”。

\subsubsection*{本文的结构}

我们勾画本文的基本论证。

在第 \ref{sec:ainf} 节中，我们定义 $A_\infty$-范畴，以及 $A_\infty$-函子、三角 $A_\infty$-范畴与分裂闭的 $A_\infty$-范畴（split-closed $A_\infty$-category）这几个重要概念。我们陈述三角包络与分裂闭包络的一般存在性结果。所有这些都是准备，为的是陈述我们将要使用的同调代数主要定理：

\begin{theorem}\label{thm:intro-main}
设 $\mathcal{A}$ 与 $\mathcal{B}$ 是两个分裂闭的 $A_\infty$-范畴，分别由 $X_1,\dots,X_n$ 与 $Y_1,\dots,Y_m$ 分裂生成（split-generate）。令
\[
  \mathcal{A}' = \Hom_{\mathcal{A}}(X_1\oplus\cdots\oplus X_n,\ X_1\oplus\cdots\oplus X_n),
\]
并对 $\mathcal{B}'$ 作同样的定义。它们是 $A_\infty$-代数。

则任何拟同构（即在上同调中诱导同构的态射）$\mathcal{A}'\to\mathcal{B}'$ 都典范地延拓为拟等价 $\mathcal{A}\to\mathcal{B}$（即在上同调中诱导等价的 $A_\infty$-函子）。
\end{theorem}

% [Tutor Note: 该定理的实质是线性 Morita 理论：$\mathcal{A}$ 由 $X_1\oplus\cdots\oplus X_n$ 分裂生成，故经 Yoneda 嵌入，$\mathcal{A}$ 可实现为 $\mathcal{A}'$ 上有限型投射模构成的 $A_\infty$-范畴；于是 $\mathcal{A}'\to\mathcal{B}'$ 的拟同构可延拓为拟等价。“典范地”掩盖了两处选择：延拓需要选取 $A_\infty$-拟逆与同伦微扰数据；并且 $\mathcal{A}'\to\mathcal{B}'$ 是 $A_\infty$-代数的拟同构（保持全部高阶乘法），延拓后只得到 $A_\infty$-范畴的拟等价（仅要求上同调等价）。]

第~\ref{sec:ainf}~节主要由定义与定理陈述构成，对此我们表示歉意。它提供的是贯穿本文其余部分的同调代数基础。由于篇幅所限，我们无法展开证明，也无法展示 $A_\infty$-结构的各种用途。建议读者只阅读第 \ref{subsec:ainf-basic} 小节，其余小节按需查阅。如果读者已经足够了解 $A_\infty$-范畴、能够看懂上述定理，那么第一节完全可以安全跳过，只在需要时回头查阅。

% [Tutor Note: 原文疑误：此处“第一节”（the first section）按上下文应指第 2 节（第一个背景章节），因为上述定理与 $A_\infty$-范畴都在第 2 节，而跳过引言并不构成“看懂定理”的理由。译文保留原文并标注。]

在第 \ref{sec:derived} 节中，我们首次接触几何，并构造一个有趣的 $A_\infty$-范畴。具体地说，利用内射消解（injective resolution），我们为概形 $X$ 上凝聚层的有界导出范畴构造一个 dg-增强（dg-enhancement）$D_\infty(X)$。我们证明：若 $X$ 在代数闭域上正则射影，则 $D_\infty(X)$ 由某个丰沛线丛（ample line bundle）的幂生成。若 $X$ 还是一维的（即一条光滑曲线），则可推知事实上线丛 $\mathcal{O}_X$ 与 $\mathcal{O}_X(P)$ 已经足够。

特别地，若 $X$ 是椭圆曲线，我们就得到一个“关系”的 dg-代数
\[
  \mathcal{Q} = \Hom_{D_\infty(X)}\bigl(\mathcal{O}_X\oplus\mathcal{O}_X(P),\ \mathcal{O}_X\oplus\mathcal{O}_X(P)\bigr),
\]
而其上同调代数 $Q = H(\mathcal{Q})$ 是一个有限维箭图代数（quiver algebra），容易显式确定。为了应用上面的定理，我们需要为 Fukaya 范畴找到类似的生成元。

当然，在这之前我们必须定义 Fukaya 范畴。这放在第 \ref{sec:fukaya} 节。我们先介绍一些基本的辛几何，然后描述 Seidel 的“分次辛几何”（graded symplectic geometry）；这是 Fukaya 范畴的基本构件。接着我们描述如何在 Fukaya 范畴中由横截拉格朗日对象所代表的那部分子集上构造 $A_\infty$-结构。这里必须预先向读者致歉：我们在两点上都不完整。第一，我们不处理非横截的情形——处理它会扯得太远。对本篇中的计算而言，只要能假定 Fukaya 范畴有一致的构造、并具有我们所声称的良好性质，一些临时办法（ad-hoc methods）就足以绕开这个问题。第二，即使对于横截情形，我们也没有给出某些符号约定的完整定义，也没有证明所声称的良好性质。后者是篇幅所限；前者则是因为这些符号除了用来建立那些良好性质之外，对我们别无用处。

接受对 Fukaya 范畴的这套描述之后，我们就可以继续推进镜像对称猜想的证明。利用 Fukaya 范畴的性质，容易证明：对于环面 $T$，两条生成同调群 $H_1(T)$ 的经线（meridian）$A$ 与 $B$ 也生成 Fukaya 范畴。于是我们得到“关系”的 $A_\infty$-代数
\[
  \mathcal{Q}' = \Hom_{D_{\Fuk}(T)}(A\oplus B,\ A\oplus B).
\]
为了应用上面的定理，我们需要找到拟同构 $\mathcal{Q}'\to\mathcal{Q}$。与代数一侧一样，计算 $H(\mathcal{Q}')$ 相当容易；直接比较即可看出 $H(\mathcal{Q}')\simeq Q$。

% [Tutor Note: “容易证明两条经线也生成 Fukaya 范畴”省略的论证是：先证 $T^2$ 的任意本质简单闭曲线都同痕于 $A$ 或 $B$ 经 Dehn 扭转后的像（映射类群 $\mathrm{Mod}(T^2)\cong SL_2(\ZZ)$ 的作用），再用分裂生成判据把生成性化归为 $A,B$ 两个对象。这条结论在 §7 会以辫作用的形式再次出现。]

此时 $A_\infty$ 图景的威力之一就显现出来：由所谓的“同伦微扰论”（homological perturbation theory），任给 $A_\infty$-代数 $\mathcal{A}$ 及其上同调代数 $A$，我们可以在 $A$ 上找到一个微分为零（$m_1=0$）的 $A_\infty$-结构 $(m_d)$，使得 $\mathcal{A}$ 拟同构于 $(A,(m_d))$。于是我们可以把 $\mathcal{A}$ 上的链级信息，换成 $A$ 上的高阶乘法。将之用于 $\mathcal{Q}$ 与 $\mathcal{Q}'$，最终得到 $Q$ 上的两个 $A_\infty$-结构 $(m_d)$ 与 $(m'_d)$；剩下的就是证明它们同构。

% [Tutor Note: 易错点：$\mathcal{Q}$（dg-代数，$m_1\neq 0$）与 $Q=H(\mathcal{Q})$（有限维箭图代数，$m_1=0$）不可混淆。原文随后讨论的是 $Q$ 上的 $A_\infty$-结构而非 $\mathcal{Q}$ 上的，这正是同伦微扰论把 $\mathcal{Q}$ 替换成 $(Q,(m_d))$ 的结果；这两步一旦混淆，§5–§6 的“分类 $Q$ 上的 $A_\infty$-结构”便无法读懂。]

此时我们把问题倒过来看：在第 \ref{sec:hoch} 节中，我们给出关于分次代数的 Hochschild 上同调的一些结果。特别地，我们说明如何用 Hochschild 上同调的计算，来分类分次代数上延拓给定乘法的那些 $A_\infty$-结构。

于是我们的目标变为：分类 $Q$ 上的所有 $A_\infty$-结构。Lekili 与 Perutz [11] 的关键洞察是：其中几乎所有的 $A_\infty$-结构实际上都通过上述过程来自椭圆曲线（有两个例外，来自奇异三次曲线）。我们在第 \ref{sec:mirror} 节证明这一关键结果，从而把论证串接起来。至此，Kontsevich 猜想在椭圆曲线情形的抽象证明完成。

% [Tutor Note: 这条是全文的收口：代数侧给出许多看似不同的 $A_\infty$-结构 $(m_d)$，辛侧只给出一个（环面上的多边形计数）。Lekili–Perutz 的定理说明它们几乎都被椭圆曲线的参数所穷尽，于是“分类”化为“找对应”。原文未展开两个例外；按上下文它们来自奇异三次曲线（节点型与尖点型，$j$-不变量退化），其几何与光滑椭圆曲线不同类。]

随后我们在第 \ref{sec:illus} 节给出各种推广与应用。首先，我们描述如何真正找出平坦环面对的镜像椭圆曲线。这本质上是一个多边形计数的手续，用到了我们在第 \ref{sec:mirror} 节途中证明的重构定理；它已由 Zaslow [24] 完成。我们勾勒所需的计算，并报告他的结果。在本节余下的部分，我们进一步图示镜像对应的细节。特别地，我们解释环面上标准的 $SL_2(\ZZ)$ 作用如何给出 Fukaya 范畴上的辫作用（braid action），并找出凝聚层的导出范畴上对应的辫作用。

\subsubsection*{记号与约定}

凡我们说链复形，指的是上链复形。类似地，我们所有的分次向量空间都按上同调的方式编号。我们采用 Koszul 符号约定来处理复形与分次向量空间的张量积。具体地说，若 $U,V,W,Z$ 是分次向量空间，$f:U\to V$、$g:W\to Z$ 是分次映射，则记 $f\otimes g:U\otimes W\to V\otimes Z$ 为映射
\[
  (f\otimes g)(x\otimes y) = (-1)^{|g||x|}f(x)\otimes g(y),
\]
其中 $|x|$ 表示元素 $x$ 的次数，$|f|$ 同理表示分次映射 $f$ 的次数。对分次向量空间，$\Hom$ 指次数为零的映射。分次向量空间 $V$ 的平移记作 $V(i)$，链复形 $C$ 的平移记作 $C[i]$。

% [Tutor Note: Koszul 符号的关键：$f\otimes g$ 作用到 $x\otimes y$ 上时，$f$ 要“跨过 $g$ 去碰 $x$”，故付出 $(-1)^{|g||x|}$。这条规则与微分约定 $d(f\otimes g)=df\otimes g+(-1)^{|f|}f\otimes dg$ 绑定；若改用 $d(f\otimes g)=df\otimes g+(-1)^{|g|}f\otimes dg$，同样的计算会给出相反的符号。另一个易错点是 $V(i)$ 与 $C[i]$ 的区别：前者只平移分次向量空间的编号（$V(i)^j=V^{i+j}$），后者把复形的微分一并平移，二者不可混用。]

% [Tutor Note: 分次按上同调编号意味着微分提升次数（$d$ 的度为 $+1$），这与按同调编号的教科书（微分降度）相差一个符号迁移；本文此后一切“次数”都按上同调编号读。在 $D^b(Coh(X))$ 上 Serre 函子为 $S=-\otimes\omega_X[n]$，椭圆曲线有 $n=1$、$\omega_X\simeq\mathcal{O}_X$，故 $S=[1]$；这正是 Verdier（Serre）对偶在平移记号上的表现，也是 A 侧 $[1]$ 平移与之匹配的原因。]

若 $X$ 是概形、$P$ 是一个点，我们记 $\mathcal{O}_{X/P}$ 为 $P$ 处的摩天层（skyscraper sheaf）$k(P)$，其中 $k(P)=\mathcal{O}_{X,P}/m_{X,P}$ 是 $P$ 处的剩余域。

% [Tutor Note: 易错点：原文用 $\mathcal{O}_{X/P}$ 记 $P$ 处的摩天层，切勿与商概形 $X/P$ 的结构层混淆（后者在本文中不出现）。同处的 $m_{X,P}$ 是极大理想的斜体写法，许多文献写作 $\mathfrak{m}_{X,P}$，两者是同一对象；剩余域因此是 $\mathcal{O}_{X,P}/m_{X,P}\cong k$。]

对于各种说得通的对象，我们用 $X^\vee$ 表示对偶。例如概形上的 $\mathcal{O}$-模，或向量空间。类似地，在种种场合我们用 $\mathbf{1}$ 或 $\mathbf{1}_X$ 表示（$X$ 上的）恒等映射。同构（或范畴的等价）记作 $\simeq$ 或 $\xrightarrow{\sim}$。对拟同构、同伦等价之类，我们不使用专门的记号。

我们用 $\Lambda_t$ 表示复数域上的万有 Novikov 域（universal Novikov field），见第 \ref{sec:fukaya} 节中关于 Novikov 域的讨论（原文第 4.4 小节）。

\subsubsection*{致谢}

许多人在我学习这些材料时给予过帮助。我特别感谢 Ivan Smith 指导本文，Ian Grojnowski 就同调代数与代数几何作过许多解释，Yanki Lekili 在讨论班上介绍他的论文并回答了我关于辛几何的问题，以及 Emile Bouaziz 的一般性讨论与文献指引。
